\documentclass[12pt,a4paper]{article}

\usepackage[utf8]{inputenc}
\usepackage[T1]{fontenc}
\usepackage{geometry}
\usepackage{setspace}
\usepackage{parskip}
\usepackage{titlesec}
\usepackage{hyperref}
\usepackage{xurl}
\usepackage{graphicx}
\usepackage{mathptmx}
\usepackage{fancyhdr}
\usepackage{xcolor}
\usepackage{tabularx}
\usepackage{booktabs}
\usepackage{float}
\usepackage{amsmath}
\usepackage{listings}
\usepackage{needspace}

\geometry{left=1.5in, right=1in, top=1in, bottom=1in}
\onehalfspacing

\renewcommand{\thesection}{\arabic{section}}
\titleformat{\section}{\fontsize{16pt}{19pt}\selectfont\bfseries}{\thesection.}{0.5em}{}
\titleformat{\subsection}{\fontsize{14pt}{17pt}\selectfont\bfseries}{\thesubsection}{0.5em}{}
\titleformat{\subsubsection}{\fontsize{12pt}{15pt}\selectfont\bfseries}{\thesubsubsection}{0.5em}{}

% ---- Code listing style ----
\definecolor{codebg}{RGB}{248,248,248}
\definecolor{codegreen}{RGB}{40,120,40}
\definecolor{codepurple}{RGB}{130,0,130}
\definecolor{codeblue}{RGB}{0,0,180}
\definecolor{codegray}{RGB}{120,120,120}

\lstdefinestyle{pythonstyle}{
    backgroundcolor=\color{codebg},
    commentstyle=\itshape\color{codegreen},
    keywordstyle=\bfseries\color{codeblue},
    stringstyle=\color{codepurple},
    basicstyle=\ttfamily\footnotesize,
    breaklines=true,
    captionpos=b,
    keepspaces=true,
    showstringspaces=false,
    language=Python,
    frame=single,
    rulecolor=\color{codegray},
    numbers=left,
    numberstyle=\tiny\color{codegray},
    numbersep=6pt,
    xleftmargin=20pt,
    framexleftmargin=20pt,
    tabsize=4,
}

\lstdefinestyle{sqlstyle}{
    backgroundcolor=\color{codebg},
    keywordstyle=\bfseries\color{codeblue},
    basicstyle=\ttfamily\footnotesize,
    breaklines=true,
    captionpos=b,
    keepspaces=true,
    showstringspaces=false,
    language=SQL,
    frame=single,
    rulecolor=\color{codegray},
    numbers=left,
    numberstyle=\tiny\color{codegray},
    numbersep=6pt,
    xleftmargin=20pt,
    framexleftmargin=20pt,
    tabsize=4,
    morekeywords={AUTOINCREMENT,REFERENCES,REAL,DATETIME,DEFAULT},
}

\lstset{style=pythonstyle}

% Placeholder box
\newcommand{\placeholder}[1]{%
    \vspace{0.5em}
    \noindent\fbox{\parbox{\linewidth}{%
        \centering\vspace{0.5em}
        \textbf{[PLACEHOLDER]} \\ \vspace{0.3em}
        \textit{#1}
        \vspace{0.5em}
    }}
    \vspace{0.5em}
}

% =====================================================================
\begin{document}
% =====================================================================

% ---- Cover Page ----
\begin{titlepage}
    \centering
    \begin{minipage}{0.48\textwidth}
        \raggedright
        \includegraphics[width=4cm]{nibm-logo.png}
    \end{minipage}
    \hfill
    \begin{minipage}{0.48\textwidth}
        \raggedleft
        \includegraphics[width=4.5cm]{coventry-logo.png}
    \end{minipage}

    \vspace{2cm}

    {\fontsize{22pt}{26pt}\selectfont\bfseries BSc Hons in Computing \par}
    {\fontsize{14pt}{17pt}\selectfont\bfseries 2024.2 \par}

    \vspace{1.5cm}

    {\fontsize{18pt}{22pt}\selectfont\bfseries
    Problem Solving with Data Structures and Algorithms (PDSA) \par}

    \vspace{0.5cm}

    {\fontsize{12pt}{16pt}\selectfont
    Individual Report: Minimum Cost Assignment Problem \\
    Coursework 1 \par}

    \vspace{3cm}

    {\large
    \begin{tabular}{lcl}
        \textbf{Coventry Index}       & \textbf{:} & \textbf{16110762} \\
        \textbf{NIBM Registered Name} & \textbf{:} & \textbf{E M A S B EKANAYAKA}
    \end{tabular}
    }

    \vfill

    {\bfseries Faculty of Engineering, Environment and Computing} \\
    {\itshape School of Computing, Electronics and Mathematics} \\
    {\bfseries Coventry University}

    \vspace{0.8cm}

    {\bfseries School of Computing and Engineering} \\
    {\itshape National Institute of Business Management}

    \vspace{0.5cm}
\end{titlepage}

\newpage
\pagenumbering{roman}
\cfoot{\thepage}

\listoffigures
\newpage

\listoftables
\newpage

\tableofcontents
\newpage

\pagenumbering{arabic}
\pagestyle{fancy}
\fancyhf{}
\renewcommand{\headrulewidth}{0pt}
\fancyfoot[R]{\thepage}

% =====================================================================
\section*{Introduction}
\addcontentsline{toc}{section}{Introduction}
% =====================================================================

This report presents an individual analysis of the Minimum Cost Assignment game
implemented as part of the PDSA coursework. The application is a full-stack web
game built with a Python FastAPI backend, a React frontend, and an SQLite database.

The core problem is the \textbf{Assignment Problem}: given $N$ employees and $N$ tasks,
where assigning employee $i$ to task $j$ incurs a cost $c_{ij}$, find a one-to-one
assignment that minimises the total cost. In each game round, $N$ is chosen randomly
in the range $[50, 100]$ and each cost $c_{ij}$ is drawn randomly from $[\$20, \$200]$.

Two algorithms are implemented and compared: a \textbf{Greedy} heuristic and the
\textbf{Hungarian algorithm} (Kuhn--Munkres method).

% =====================================================================
\clearpage
\section{Program Logic}
% =====================================================================

\subsection{Game Round Generation}

Each round begins when the frontend calls the \texttt{GET /api/round/new} endpoint.
The backend uses Python's \texttt{random} module to produce a random problem instance:

\needspace{7cm}
\begin{lstlisting}[caption={Game round generation -- backend/main.py}]
@app.get("/api/round/new", response_model=NewRoundResponse)
def new_round():
    n = random.randint(50, 100)
    cost_matrix = [
        [random.randint(20, 200) for _ in range(n)]
        for _ in range(n)
    ]
    return NewRoundResponse(n=n, cost_matrix=cost_matrix)
\end{lstlisting}

\noindent
The $N \times N$ cost matrix is sent to the frontend, which renders it as an
interactive table. The player may optionally assign some tasks manually before
requesting the algorithmic solutions.

\subsection{Greedy Assignment Algorithm}

The greedy approach scans employees in order $0, 1, \ldots, N-1$. For each employee
it picks the \emph{cheapest unassigned task} and commits to that choice immediately.
Because earlier decisions are never reconsidered, the greedy algorithm does not
guarantee a globally optimal assignment; it serves as a fast baseline.

\needspace{14cm}
\begin{lstlisting}[caption={Greedy core loop -- backend/algorithms/greedy.py}]
for employee in range(n):
    best_task = -1
    best_cost = float("inf")
    candidates = []

    for task in range(n):
        if task not in assigned_tasks:
            candidates.append((task, cost_matrix[employee][task]))
            if cost_matrix[employee][task] < best_cost:
                best_cost = cost_matrix[employee][task]
                best_task = task

    assignment[employee] = best_task
    assigned_tasks.add(best_task)
    running_total += int(best_cost)

    steps.append({
        "employee":      employee,
        "task":          best_task,
        "cost":          int(best_cost),
        "running_total": running_total,
        "candidates":    sorted(candidates, key=lambda x: x[1])[:5],
    })

total_cost = sum(cost_matrix[emp][assignment[emp]] for emp in range(n))
return assignment, total_cost, steps
\end{lstlisting}

\noindent
The function returns three values:
\begin{itemize}
    \item \texttt{assignment} -- a list where \texttt{assignment[i]} is the task
          index allocated to employee $i$.
    \item \texttt{total\_cost} -- the sum of all assignment costs in dollars.
    \item \texttt{steps} -- a list of dictionaries (one per employee) that the
          frontend uses to animate the decision process step by step.
\end{itemize}

\subsection{Hungarian Algorithm}

The Hungarian algorithm (Kuhn--Munkres, 1955) is guaranteed to find the
\emph{globally optimal} assignment. It operates on a reduced cost matrix and
uses König's theorem to determine when an optimal assignment has been found.

The algorithm proceeds through five logical phases that repeat until an optimal
solution is identified:

\needspace{8cm}
\begin{enumerate}
    \item \textbf{Row reduction.} Subtract the minimum value of each row from every
          cell in that row, ensuring at least one zero per row.
    \item \textbf{Column reduction.} Subtract the minimum value of each column from
          every cell in that column, ensuring at least one zero per column.
    \item \textbf{Minimum line cover.} Find the minimum number of horizontal or
          vertical lines that cover all zeros. This is equivalent to computing a
          maximum bipartite matching on the zero graph via augmenting paths
          (König's theorem).
    \item \textbf{Optimality check.} If the number of covering lines equals $N$,
          a perfect assignment through zeros exists and is optimal.
    \item \textbf{Matrix adjustment.} If the cover count is less than $N$, find the
          minimum uncovered value $\delta$, subtract $\delta$ from all uncovered cells,
          add $\delta$ to doubly-covered cells, and repeat from step 3.
\end{enumerate}

\needspace{12cm}
\begin{lstlisting}[caption={Hungarian main body -- backend/algorithms/hungarian.py}]
n = len(cost_matrix)
matrix = [row[:] for row in cost_matrix]

# Step 1 -- Row reduction
for i in range(n):
    row_min = min(matrix[i])
    for j in range(n):
        matrix[i][j] -= row_min

# Step 2 -- Column reduction
for j in range(n):
    col_min = min(matrix[i][j] for i in range(n))
    for i in range(n):
        matrix[i][j] -= col_min

# Iterate until minimum line cover equals n
while True:
    row_cover, col_cover = _min_line_cover(matrix, n)
    if len(row_cover) + len(col_cover) >= n:
        break
    _adjust_matrix(matrix, n, row_cover, col_cover)

assignment = _extract_assignment(matrix, n)
total_cost = sum(cost_matrix[emp][assignment[emp]] for emp in range(n))
\end{lstlisting}

\noindent
The helper \texttt{\_min\_line\_cover} finds the maximum matching on zero entries
using an iterative depth-first search for augmenting paths, then applies König's
theorem to derive the minimum line cover. Once the loop terminates,
\texttt{\_extract\_assignment} reads the zero positions of the final reduced matrix
to produce the bijective assignment.

\subsection{Timing Measurement}

Each algorithm is timed individually using Python's high-resolution
\texttt{time.perf\_counter()} clock. Only the algorithm execution itself is measured,
excluding network I/O, parsing, and database writes.

\needspace{6cm}
\begin{lstlisting}[caption={Timing capture -- backend/main.py}]
for algo_name, algo_fn in [("Greedy", greedy_assignment),
                            ("Hungarian", hungarian_assignment)]:
    sub_matrix = [[matrix[i][j] for j in free_tasks]
                  for i in free_employees]
    t0 = time.perf_counter()
    sub_assignment, _, sub_steps = algo_fn(sub_matrix)
    elapsed_ms = (time.perf_counter() - t0) * 1000
\end{lstlisting}

\noindent
The elapsed time is rounded to four decimal places and stored in the database as a
\texttt{REAL} value in milliseconds.

\subsection{Database Persistence}

After both algorithms complete, results are stored in two normalised SQLite tables:

\needspace{8cm}
\begin{lstlisting}[style=sqlstyle, caption={Database schema -- backend/database.py}]
CREATE TABLE IF NOT EXISTS game_rounds (
    id              INTEGER PRIMARY KEY AUTOINCREMENT,
    n               INTEGER  NOT NULL,
    player_name     TEXT,
    user_total_cost INTEGER,
    created_at      DATETIME DEFAULT CURRENT_TIMESTAMP
);

CREATE TABLE IF NOT EXISTS algorithm_results (
    id             INTEGER PRIMARY KEY AUTOINCREMENT,
    round_id       INTEGER NOT NULL REFERENCES game_rounds(id),
    algorithm_name TEXT    NOT NULL,
    total_cost     INTEGER NOT NULL,
    time_ms        REAL    NOT NULL
);
\end{lstlisting}

\noindent
Each round produces one row in \texttt{game\_rounds} and two rows in
\texttt{algorithm\_results} (one per algorithm). The foreign key
\texttt{round\_id} links timing records to their parent round.

% =====================================================================
\clearpage
\section{Complexity Analysis}
% =====================================================================

\subsection{Greedy Algorithm -- $O(n^2)$ Time, $O(n)$ Space}

\subsubsection{Time Complexity}

The greedy algorithm contains two nested loops. The outer loop iterates over $n$
employees and the inner loop scans all $n$ tasks to find the cheapest available
option. No further passes are required:
\[
    T_{\text{Greedy}}(n) = O(n) \times O(n) = O(n^2)
\]
For $n = 100$ this amounts to at most $10{,}000$ comparisons, completing in
microseconds on modern hardware.

\subsubsection{Space Complexity}

The only auxiliary structure is the \texttt{assigned\_tasks} set, which holds at
most $n$ integers. The \texttt{steps} list adds $n$ small constant-size dictionaries.
Both are $O(n)$, giving an overall space complexity of $O(n)$.

\subsubsection{Optimality}

The greedy algorithm is a \textbf{heuristic}. Committing to each decision without
lookahead can miss the global optimum. For instance, selecting the cheapest task
for employee 0 may force a more expensive assignment for employee 1. The unit test
\texttt{test\_hungarian\_at\_least\_as\_good\_as\_greedy} (20 random matrices,
seed 42) confirms empirically that the greedy cost is always greater than or equal
to the Hungarian optimal.

\subsection{Hungarian Algorithm -- $O(n^3)$ Time, $O(n^2)$ Space}

\subsubsection{Time Complexity}

\begin{itemize}
    \item Row and column reduction: two passes over the $n \times n$ matrix -- $O(n^2)$.
    \item Maximum matching via augmenting paths: iterative DFS over at most $n$
          rows and $n$ columns per employee -- $O(n^2)$ per matching, $O(n^3)$
          across all employees.
    \item König's theorem propagation: BFS/DFS over the zero graph -- $O(n^2)$.
    \item Matrix adjustment: one full pass over $n^2$ cells -- $O(n^2)$.
    \item Number of iterations: each iteration increases the matching size by at
          least one, so there are at most $n$ outer iterations.
\end{itemize}

The dominant term is $n$ iterations of an $O(n^2)$ matching:
\[
    T_{\text{Hungarian}}(n) = O(n^3)
\]
For $n = 100$ this is at most $1{,}000{,}000$ basic operations. Observed runtimes
are typically in the range of a few milliseconds.

\subsubsection{Space Complexity}

The algorithm maintains a working copy of the $n \times n$ matrix throughout,
giving $O(n^2)$ space. Auxiliary arrays for row/column matching and cover sets
add $O(n)$ each, dominated by the matrix copy. Overall space complexity: $O(n^2)$.

\subsubsection{Optimality}

The Hungarian algorithm is \textbf{provably optimal}. The termination condition --
a minimum line cover of size $n$, equivalent to a perfect matching on zeros of the
reduced matrix -- guarantees that the extracted assignment achieves the minimum
possible total cost.

\subsection{Analysis Based on Program Outputs}

Table \ref{tab:complexity} summarises the theoretical and observed properties of
both algorithms as experienced during game play.

\begin{table}[H]
    \centering
    \caption{Complexity summary from program logic and observed outputs}
    \label{tab:complexity}
    \begin{tabularx}{\linewidth}{lXX}
        \toprule
        \textbf{Property}          & \textbf{Greedy}      & \textbf{Hungarian}   \\
        \midrule
        Time complexity            & $O(n^2)$             & $O(n^3)$             \\
        Space complexity           & $O(n)$               & $O(n^2)$             \\
        Typical runtime ($n=100$)  & $< 0.5$\,ms          & $2$--$8$\,ms         \\
        Optimality                 & Heuristic only       & Guaranteed optimal   \\
        Cost vs.\ optimal          & Always $\geq$ optimal & Always optimal      \\
        Growth with $n$            & Quadratic            & Cubic                \\
        \bottomrule
    \end{tabularx}
\end{table}

% =====================================================================
\clearpage
\section{Comparison of Algorithmic Approaches}
% =====================================================================

\subsection{Approach and Design Philosophy}

The two algorithms represent contrasting design philosophies:

\begin{itemize}
    \item The \textbf{Greedy} algorithm applies the principle of \emph{local optimality}:
          make the best available choice at each step without reconsidering prior
          decisions. This leads to fast, simple code but sacrifices global optimality.
    \item The \textbf{Hungarian} algorithm applies \emph{global optimisation} through
          algebraic transformations (matrix reduction) and combinatorial reasoning
          (König's theorem and maximum bipartite matching). It is slower but
          always correct.
\end{itemize}

\subsection{Cost Quality}

The unit test \texttt{test\_hungarian\_at\_least\_as\_good\_as\_greedy}, run over
20 randomly generated matrices with $n \in [3, 15]$ using seed 42, confirms:
\[
    \text{cost}_{\text{Hungarian}} \leq \text{cost}_{\text{Greedy}} \quad
    \text{for all tested instances}
\]
In practice the greedy solution is often 5--20\% more expensive than the optimum
for random cost matrices. The gap tends to be larger when the matrix has few very
low-cost cells, forcing greedy into costly early choices.

\subsection{Runtime Performance}

For the game's problem size ($n \in [50, 100]$) both algorithms are fast. Greedy
completes in under 0.5\,ms while Hungarian typically takes 2--8\,ms. At larger
scales the cubic growth of Hungarian becomes significant: scaling from $n = 100$
to $n = 200$ increases Greedy's runtime by $4\times$ but Hungarian's by $8\times$.

\subsection{Summary Comparison}

\begin{table}[H]
    \centering
    \caption{Comparison of Greedy and Hungarian algorithms}
    \label{tab:comparison}
    \begin{tabularx}{\linewidth}{lXX}
        \toprule
        \textbf{Criterion}          & \textbf{Greedy}                     & \textbf{Hungarian}                   \\
        \midrule
        Time complexity             & $O(n^2)$                            & $O(n^3)$                             \\
        Space complexity            & $O(n)$                              & $O(n^2)$                             \\
        Optimality                  & No (heuristic)                      & Yes (provably optimal)               \\
        Implementation complexity   & Simple (two nested loops)           & Complex (matching + König's theorem) \\
        Scalability                 & Excellent for large $n$             & Suitable up to $n \approx 1{,}000$   \\
        Predictability              & Constant number of passes           & Iterations vary per matrix           \\
        Typical runtime at $n=100$  & $< 0.5$\,ms                         & $2$--$8$\,ms                         \\
        Use case                    & Real-time scheduling, large $n$     & Exact optimisation, moderate $n$     \\
        \bottomrule
    \end{tabularx}
\end{table}

\subsection{When to Prefer Each Algorithm}

\begin{itemize}
    \item \textbf{Prefer Greedy} when speed is critical and a near-optimal solution
          is acceptable -- for example, in real-time systems, large-scale logistics,
          or as a warm-start heuristic for more complex solvers.
    \item \textbf{Prefer Hungarian} when the exact minimum cost is required and $n$
          is not prohibitively large -- for example, in task scheduling, exam
          timetabling, or any domain where a suboptimal assignment incurs a
          measurable financial or operational penalty.
\end{itemize}

In the context of this game, the Hungarian result is presented as the optimal
reference solution, while Greedy provides a fast comparison baseline that
illustrates the cost of approximation.

% =====================================================================
\clearpage
\section{Timing Chart -- Game Rounds}
\label{sec:chart}
% =====================================================================

The chart in Figure~\ref{fig:timing} shows the execution time (in milliseconds)
of each algorithm across the rounds played individually. Greedy times are drawn
in amber and Hungarian times in blue. Data were retrieved from the
\texttt{algorithm\_results} table in the SQLite database.

\begin{figure}[H]
    \centering
    \includegraphics[width=\linewidth]{algorithm-timing.png}
    \caption{Algorithm timing across game rounds (Greedy vs.\ Hungarian).}
    \label{fig:timing}
\end{figure}

\noindent
\textbf{Observations from the chart:}
\begin{itemize}
    \item Greedy times are so small that the amber bars are barely visible on the
          shared y-axis. Measured values are consistently below 1.5\,ms and often
          under 0.5\,ms.
    \item Hungarian times dominate the chart, typically ranging from 4 to 12\,ms
          with an average around 8--9\,ms.
    \item Hungarian is consistently 20--50 times slower than Greedy, which matches
          the theoretical $O(n^3)$ vs.\ $O(n^2)$ difference.
    \item Round~6 peaks at roughly 17\,ms for Hungarian -- the worst case in this
          sample -- which reflects a cost matrix that required more iterations of
          the minimum line cover loop.
    \item Round~9 shows an unusually low Hungarian time (around 2\,ms), indicating
          a matrix whose zeros happened to form a perfect matching early in the
          reduction phase.
    \item Rounds~1, 7, and~14 show negligible bars for both algorithms. These
          correspond to rounds in which the player assigned all (or most) employees
          manually, leaving only a tiny sub-matrix for the algorithms to solve.
\end{itemize}

% =====================================================================
\clearpage
\section{Video Demonstration}
\label{sec:video}
% =====================================================================

The video clip demonstrates the complete game rounds played for this submission,
showing the algorithm animation, step-by-step assignment, timing results, and
database update for each round.

\vspace{0.5em}
\noindent
\textbf{Video link:} \\
\url{https://nibm-my.sharepoint.com/:f:/g/personal/cobsccomp242p-065_student_nibm_lk/IgC1z1sqJ0tgQ5aY-_paQP58ASlJvXpEXgEXOfWhQnmc9rQ?e=xGp5Ho}
\vspace{0.5em}

\noindent
The video covers:
\begin{itemize}
    \item Starting a new round (random $N$ and cost matrix generation).
    \item Running both algorithms and observing the animated step-by-step assignment.
    \item Reviewing the comparison panel (cost, time, and gap analysis).
    \item The timing chart updating after each round.
    \item The database row created for each round.
\end{itemize}

% =====================================================================
\clearpage
\section{Database Output}
\label{sec:db}
% =====================================================================

Figure~\ref{fig:db} shows the content of the \texttt{algorithm\_results} table,
confirming that both algorithm names, total costs, and execution times are
correctly persisted for every round.

\begin{figure}[H]
    \centering
    \includegraphics[width=0.95\linewidth]{algorithm-results.png}
    \caption{Rows of the \texttt{algorithm\_results} table produced during game play.}
    \label{fig:db}
\end{figure}

\noindent
\textbf{Observations from the database:}
\begin{itemize}
    \item Each round produces exactly two rows linked by \texttt{round\_id} --
          one for Greedy and one for Hungarian -- giving 38 rows across 19 rounds.
    \item \texttt{Hungarian.total\_cost} is always less than or equal to
          \texttt{Greedy.total\_cost}, empirically confirming Hungarian's optimality.
    \item The cost saving varies per round. Round~9 shows the largest gap
          (Greedy \$2274 vs.\ Hungarian \$1539, a 32\% saving), while Round~6 shows
          the smallest (Greedy \$3019 vs.\ Hungarian \$2761, an 8.5\% saving).
    \item Rows for Rounds~1 and~7 show very small values (e.g., cost 9--10 for
          Round~1, around 100 for Round~7). These correspond to rounds where the
          player manually assigned most employees, leaving only a small sub-matrix
          for the algorithm.
    \item \texttt{time\_ms} is stored with four-decimal-place precision as a
          \texttt{REAL} value, preserving sub-millisecond Greedy timings (e.g.,
          0.0075\,ms, 0.2377\,ms).
\end{itemize}

\noindent
Table \ref{tab:db-schema} describes the columns visible in the screenshot.

\begin{table}[H]
    \centering
    \caption{Column descriptions for the \texttt{algorithm\_results} table}
    \label{tab:db-schema}
    \begin{tabularx}{\linewidth}{llX}
        \toprule
        \textbf{Column}          & \textbf{Type}    & \textbf{Description}                             \\
        \midrule
        \texttt{id}              & INTEGER          & Auto-incremented primary key                     \\
        \texttt{round\_id}       & INTEGER (FK)     & Foreign key to \texttt{game\_rounds.id}          \\
        \texttt{algorithm\_name} & TEXT             & ``Greedy'' or ``Hungarian''                      \\
        \texttt{total\_cost}     & INTEGER          & Total assignment cost in dollars                 \\
        \texttt{time\_ms}        & REAL             & Execution time in milliseconds (4 d.p.)          \\
        \bottomrule
    \end{tabularx}
\end{table}

% =====================================================================
\clearpage
\addcontentsline{toc}{section}{List of References}
\section*{List of References}
% =====================================================================

\begin{enumerate}
    \item Kuhn, H. W. (1955). The Hungarian method for the assignment problem.
          \textit{Naval Research Logistics Quarterly}, 2(1--2), 83--97.
    \item Munkres, J. (1957). Algorithms for the assignment and transportation
          problems. \textit{Journal of the Society for Industrial and Applied
          Mathematics}, 5(1), 32--38.
    \item Cormen, T. H., Leiserson, C. E., Rivest, R. L., and Stein, C. (2009).
          \textit{Introduction to Algorithms} (3rd ed.). MIT Press.
    \item Konig, D. (1931). Graphen und Matrizen.
          \textit{Matematikai Lapok}, 38, 116--119.
\end{enumerate}

\end{document}
